\clearpage

\phantomsection

\setcounter{chapter}{4}
\chapter[{KẾT LUẬN VÀ HƯỚNG PHÁT TRIỂN}]{KẾT LUẬN VÀ HƯỚNG PHÁT TRIỂN}

Chương này tổng kết các kết quả đạt được của luận văn, phân tích những hạn chế còn tồn tại và đề xuất hướng phát triển trong tương lai.

% =====================================================
\section{Kết luận}
\label{sec:conclusion}

Luận văn đã nghiên cứu, triển khai và đánh giá một hệ thống chuyển đổi ảnh chân dung sang phong cách anime dựa trên mô hình DTGAN (AnimeGANv3) \cite{Liu2024}, kết hợp với pipeline tích hợp phát hiện khuôn mặt sử dụng RetinaFace \cite{Deng2019RetinaFace} hỗ trợ ba chế độ vận hành (Face Mode, Landscape Mode, Hybrid Mode). Qua quá trình thực hiện, luận văn đạt được các kết quả chính sau:

\begin{enumerate}
    \item \textbf{Phân tích kiến trúc DTGAN:} Trình bày hệ thống và chi tiết kiến trúc Double-Tail Generator (Base Encoder, Support Tail, Main Tail), kỹ thuật chuẩn hóa LADE, cơ chế External Attention v3 và hệ thống 7 hàm mất mát chuyên biệt cho bài toán anime hóa. Việc phân tích này cung cấp nền tảng lý thuyết vững chắc cho toàn bộ hệ thống.

    \item \textbf{Xây dựng pipeline tích hợp phát hiện khuôn mặt:} Đây là đóng góp nguyên bản của luận văn — tích hợp RetinaFace MobileNet 0.25 qua ONNX Runtime với thuật toán Margin Expansion 6 bước, hỗ trợ ba chế độ: Face Mode (xử lý khuôn mặt lớn nhất, FID cải thiện $8.7\%$ so với Landscape Mode), Landscape Mode (xử lý toàn ảnh), và Hybrid Mode (kết hợp cả hai bằng kỹ thuật feathered blending cho ảnh nhóm). Pipeline cung cấp cơ chế graceful degradation tự động chuyển sang Landscape Mode khi không phát hiện được khuôn mặt.

    \item \textbf{Huấn luyện và đánh giá toàn diện:} Quá trình huấn luyện qua 2 giai đoạn (5 epoch pre-training + 69 epoch adversarial) cho thấy hội tụ ổn định: Pre-train G\_loss giảm từ $1.74$ xuống $0.25$; G\_loss ổn định ở $\approx 6.5$ và D\_loss ở $\approx 0.11$ — không xảy ra mode collapse. Đánh giá định lượng cho kết quả FID $= 58.6$ và LPIPS $= 0.384$ (Face Mode), vượt trội so với CycleGAN, CartoonGAN và AnimeGANv2.

    \item \textbf{Triển khai ứng dụng thực tế:} Xây dựng ứng dụng web tích hợp (React frontend + FastAPI backend), hỗ trợ chuyển đổi cả ảnh và video với ba chế độ Face Mode, Landscape Mode và Hybrid Mode. Model ONNX chỉ 9.1 MB, đạt tốc độ 42 FPS tại $256 \times 256$ trên GPU, phù hợp cho triển khai thực tế.
\end{enumerate}

\subsection{Mức độ hoàn thành mục tiêu}
\label{subsec:goal_completion}

Bảng~\ref{tab:goal_completion} so sánh mức độ hoàn thành với từng mục tiêu đặt ra trong phần Mở đầu.

\begin{table}[H]
\caption{Đánh giá mức độ hoàn thành các mục tiêu nghiên cứu}
\label{tab:goal_completion}
\centering
\begin{tabular}{|p{6.5cm}|c|p{5.5cm}|}
\hline
\textbf{Mục tiêu} & \textbf{Kết quả} & \textbf{Minh chứng} \\
\hline
Nghiên cứu nền tảng lý thuyết CNN, GAN, RetinaFace & \checkmark\ Hoàn thành & Chương 1: trình bày đầy đủ \\
\hline
Phân tích kiến trúc DTGAN & \checkmark\ Hoàn thành & Chương 2: Generator, Discriminator, LADE, EA v3, 7 loss \\
\hline
Xây dựng pipeline Face Extraction & \checkmark\ Hoàn thành & Chương 3: RetinaFace + Margin Expansion + Hybrid Mode, FID cải thiện $8.7\%$ \\
\hline
Huấn luyện và đánh giá mô hình & \checkmark\ Hoàn thành & Chương 4: FID 58.6, LPIPS 0.384, 42 FPS \\
\hline
Xây dựng ứng dụng demo & \checkmark\ Hoàn thành & Ứng dụng web React + FastAPI, hỗ trợ ảnh và video \\
\hline
\end{tabular}
\end{table}

% =====================================================
\section{Hạn chế}
\label{sec:limitations}

Mặc dù đạt được các kết quả khả quan, hệ thống vẫn còn một số hạn chế cần được nhận diện. Thứ nhất, phong cách anime học được phụ thuộc hoàn toàn vào tập dữ liệu style; nếu tập dữ liệu không đủ đa dạng về góc chiếu, biểu cảm và điều kiện ánh sáng, mô hình sẽ có xu hướng overfitting vào một phong cách hẹp. Thứ hai, RetinaFace MobileNet 0.25 giảm hiệu năng đáng kể tại góc nghiêng $>60°$ (recall $<76\%$), dẫn đến bounding box không chính xác và chất lượng chuyển đổi kém ở những trường hợp này. Thứ ba, pipeline xử lý video hiện tại chuyển đổi từng frame độc lập, chưa có cơ chế đảm bảo tính nhất quán giữa các frame liên tiếp (temporal consistency), có thể gây hiệu ứng nhấp nháy (flickering) khi xem video đầu ra. Thứ tư, chi phí tính toán của chế độ Hybrid Mode tăng tuyến tính theo số khuôn mặt ($1 + N$ lần suy luận ONNX), có thể chậm với ảnh nhóm đông ($N > 5$). Cuối cùng, tốc độ xử lý trên CPU chỉ đạt $\approx 1.1$ FPS ở $256\times256$, chưa phù hợp cho người dùng không có GPU; model ONNX chưa được áp dụng quantization để tối ưu cho CPU inference.

% =====================================================
\section{Hướng phát triển}
\label{sec:future_work}

Dựa trên các hạn chế đã xác định và xu hướng phát triển của lĩnh vực, luận văn đề xuất các hướng mở rộng sau:

\subsection{Cải thiện chất lượng xử lý khuôn mặt}
\label{subsec:future_face}

\textbf{Face Alignment:} RetinaFace đã trả về 5 facial landmarks (2 mắt, mũi, 2 khóe miệng) nhưng thông tin này chưa được khai thác. Sử dụng các landmarks này để căn chỉnh khuôn mặt về góc nhìn trực diện chuẩn trước khi đưa vào Generator sẽ cải thiện đáng kể chất lượng đầu ra cho khuôn mặt nghiêng:

\begin{equation}
\label{eq:face_align}
I_{\text{aligned}} = \mathcal{T}_{\text{affine}}\!\left(I_{\text{crop}},\; \text{landmarks},\; \text{target\_template}\right)
\end{equation}

\textbf{Nâng cấp Face Blending:} Chế độ Hybrid Mode hiện tại đã giải quyết vấn đề ghép khuôn mặt vào nền bằng feathered blending (alpha mask + Gaussian blur). Tuy nhiên, có thể cải thiện chất lượng ghép bằng Poisson Image Editing \cite{Perez2003Poisson} -- kỹ thuật seamless cloning giải phương trình Poisson tại biên, cho kết quả tự nhiên hơn alpha blending đơn giản, đặc biệt khi có khác biệt lớn về ánh sáng và tông màu giữa vùng mặt và nền.

\textbf{Nâng cấp detector:} Thay backbone MobileNet 0.25 bằng ResNet50 hoặc sử dụng YOLOv8-face \cite{Jocher2023YOLOv8} -- detector khuôn mặt thế hệ mới cân bằng tốt hơn giữa tốc độ và độ chính xác, đặc biệt ở khuôn mặt nghiêng.

\subsection{Xử lý video với temporal consistency}
\label{subsec:future_video}

Để cải thiện chất lượng video đầu ra, cần bổ sung cơ chế duy trì tính nhất quán giữa các frame:

\begin{equation}
\label{eq:temporal}
I_{\text{anime}}^{(t)} = \mathcal{G}\!\left(I^{(t)},\; \mathcal{W}\!\left(I_{\text{anime}}^{(t-1)},\; \mathbf{v}^{(t)}\right)\right)
\end{equation}

trong đó $\mathbf{v}^{(t)}$ là optical flow từ frame $t-1$ sang $t$ \cite{Horn1981opticalflow}, $\mathcal{W}$ là phép warping. Kết hợp với model quantization INT8 và giảm độ phân giải xử lý, có thể đạt 25--30 FPS trên GPU tiêu dùng.

\subsection{Mở rộng đa phong cách và triển khai mobile}
\label{subsec:future_multistyle}

\textbf{Multi-style model:} Thay vì huấn luyện riêng biệt cho mỗi phong cách, xây dựng một mô hình đa phong cách bằng cách đưa one-hot vector phong cách vào Generator qua Conditional Normalization, hoặc áp dụng nội suy tuyến tính trong không gian tham số LADE:
\begin{equation}
\label{eq:style_interp}
\theta_{\text{LADE}}(\lambda) = (1-\lambda)\,\theta_{s_1} + \lambda\,\theta_{s_2}, \quad \lambda \in [0,1]
\end{equation}

\textbf{Triển khai mobile:} Model ONNX 9.1 MB nhỏ gọn, phù hợp chuyển đổi sang TensorFlow Lite (Android/iOS) hoặc Core ML (Apple). Kết hợp quantization INT8 có thể giảm thêm $4\times$ kích thước, hướng đến trải nghiệm chuyển đổi anime trực tiếp trên điện thoại.

% =====================================================
\section{Lời kết}
\label{sec:final_word}

Luận văn đã hoàn thành mục tiêu xây dựng một hệ thống chuyển đổi ảnh chân dung sang phong cách anime hoàn chỉnh, từ nghiên cứu kiến trúc mô hình DTGAN đến tích hợp pipeline phát hiện khuôn mặt với ba chế độ vận hành (Face Mode, Landscape Mode, Hybrid Mode) và triển khai thành ứng dụng web thực tế.

Kết quả định lượng (FID $= 58.6$, LPIPS $= 0.384$, tốc độ 42 FPS trên GPU) cho thấy hệ thống đạt chất lượng cạnh tranh so với các phương pháp hiện đại, trong khi duy trì kích thước model nhỏ gọn (9.1 MB ONNX) phù hợp cho triển khai thực tế. Pipeline tích hợp với thuật toán Margin Expansion và chế độ Hybrid Mode (feathered blending) là đóng góp nguyên bản, đã được xác nhận cải thiện chất lượng khoảng $8.7\%$ FID so với Landscape Mode trên ảnh chân dung, đồng thời cho phép xử lý ảnh nhóm với chất lượng khuôn mặt cao mà vẫn giữ nền.

Các hướng phát triển được đề xuất — face alignment, nâng cấp blending bằng Poisson editing, temporal consistency cho video và mở rộng đa phong cách — tạo nên lộ trình nghiên cứu rõ ràng và khả thi, hướng tới một hệ thống anime stylization toàn diện phục vụ người dùng rộng rãi.

% \vspace{1cm}
% \noindent\textbf{Lời cảm ơn/Tài trợ:}
Nghiên cứu này được tiến hành trong khuôn khổ đề tài QG.25.03 "Nghiên cứu tái tạo chuyển động diễn xuất khuôn mặt nhân vật trên mô hình 3D, phục vụ công tác giảng dạy sân khấu truyền thống".
